Let $T_n$ count all valid length-$n$ strings. Classifying by a final run of length one or two gives $T_n=T_{n-1}+T_{n-2}$, with $T_1=2,T_2=4$. Thus $T_{12}=466$.
Count balanced strings with six of each digit by runs. The number of compositions of six into $k$ parts of size one or two is $A_k=\binom{k}{6-k}$, with nonzero values $A_3=1,A_4=6,A_5=5,A_6=1$. If the first and last bits differ, the two colours have equally many runs, giving $2\sum A_k^2=2(1+36+25+1)=126$. If they agree, one colour has one extra run, giving $2\sum A_kA_{k+1}=2(6+30+5)=82$. Hence there are $208$ balanced strings.
Complementing every bit pairs strings with more ones with strings with more zeros. The requested count is $\boxed{\dfrac{466-208}{2}=129}$.
